\section{Survey Methodology}\label{sec:method}

\subsection{Research Objectives}

Our survey delves into the utilization of graph prompts within the context of Artificial General Intelligence (AGI), focusing on their potential to enhance the generality of graph models across various tasks, domains, and modalities. We aim to investigate how graph prompts can integrate with foundational models to interact with diverse data types such as text and images. Our objectives include:

\begin{itemize}
    \item Developing a unified framework to consolidate varied approaches in graph prompting, facilitating a comprehensive understanding of existing methodologies.
    \item Exploring the inherent nature and strategic roles of prompts in the realm of graph-based learning, aiming to enhance the interpretability and functionality of graph models.
    \item Innovating in the design of graph prompts, focusing on creating efficient, versatile, and context-sensitive prompts that cater to specific requirements of graph-based tasks.
    \item Proposing robust criteria for evaluating the effectiveness and efficiency of graph prompts, including their adaptability and scalability in different application scenarios.
    \item Examining the practical implications of graph prompts in real-world settings, assessing their implementation challenges, and highlighting successful case studies.
    \item Identifying current limitations in graph prompting and forecasting future trends and potential research directions in this evolving field.
\end{itemize}

\subsection{Taxonomy}

The taxonomy of this survey is intricately designed to categorize graph prompts based on their specific applications and functionalities, providing a structured approach to understanding their role in AGI.

\subsubsection{Prompt Design for Graph Tasks}
\begin{enumerate}
    \item \textbf{Node-level Task Graph Prompt Formulation:} This segment delves into crafting prompts that enhance understanding at the node level, aiming to extract significant insights about individual nodes' roles within complex structures.
    \item \textbf{Edge-level Task Graph Prompt Crafting:} Focuses on edge-level tasks, exploring how tailored prompts illuminate intricate relationships and dependencies in graph structures.
    \item \textbf{Graph-level Task Graph Prompt Application:} Examines prompts designed for broad graph-level analysis, emphasizing their utility in holistic understanding and interpretation of entire graph structures.
    \item \textbf{Multi-task Graph Prompting:} Investigates the application of prompts in multi-task scenarios, showcasing their adaptability in providing comprehensive insights from graph data.
    \item \textbf{Quality Assessment of Graph Prompts:} Introduces a framework for evaluating the quality of graph prompts, considering factors like efficiency, execution time, and data reformulation capabilities.
\end{enumerate}

\subsubsection{Multi-modal Prompting Utilizing Graphs (Text-Graph)}
\begin{enumerate}
    \item \textbf{Textual and Graph Data Fusion:} Explores the integration of text and graph data, assessing the enhanced analytical capabilities and decision-making insights this combination provides.
    \item \textbf{Large Language Models in Graph Data Processing:} Examines the use of LLMs in processing graph data, highlighting their potential in natural language understanding and generative applications.
\end{enumerate}

\subsubsection{Graph Domain Adaptation through Prompting Techniques}
\begin{enumerate}
    \item \textbf{Cross-Domain Graph Adaptation Strategies:} Investigates how prompting techniques facilitate knowledge transfer across various graph domains, focusing on domain-specific adaptability.
    \item \textbf{Context-aware Graph Adaptation:} Explores the tailoring of prompts for context-specific graph analysis, enhancing relevance and accuracy.
    \item \textbf{Synergy between Graph Embeddings and Domain Adaptation:} Analyzes the interplay between graph embeddings and domain adaptation, facilitated by prompts.
    \item \textbf{Transfer Learning in Graph Analysis with Prompts:} Examines how existing knowledge and models can be adapted to new domains through prompts, optimizing analysis efficiency.
    \item \textbf{Dynamic Domain Adaptation with Prompts:} Discusses real-time adaptability strategies in the evolving nature of domain adaptation.
    \item \textbf{Case Studies and Practical Applications:} Presents practical case studies demonstrating the successful application of graph domain adaptation using prompts.
    \item \textbf{Future Challenges and Directions in Graph Prompting:} Identifies ongoing challenges and potential future developments in graph prompting.
\end{enumerate}

\subsection{Literature Overview}

We conducted an exhaustive literature review, encompassing over 100 high-quality papers from prestigious conferences and journals within the past five years. These sources include SIGKDD, AAAI, ACL, CIKM, EMNLP, ICLR, ICML, IJCAI, NeurIPS, SIGIR, and TKDE, most of which are ranked as CCF A or CORA A*. The reviewed literature covers a broad spectrum of topics relevant to graph prompts in AGI, including theoretical foundations, algorithmic innovations, application scenarios, and empirical studies.

\subsubsection{Connection to Existing Work}
Our survey identifies and addresses gaps in the existing literature, distinguishing our work from previous surveys by focusing on the unique intersection of graph prompts and AGI. We highlight our novel contributions, including the development of a comprehensive taxonomy, an extensive review of cutting-edge methodologies, and the identification of future research trajectories in the realm of graph prompting.

\subsubsection{Thematic Synthesis}
The reviewed literature is systematically categorized into thematic areas, aligning with our research objectives and taxonomy. Each theme is critically analyzed to distill key findings, methodologies, and insights, contributing to a nuanced understanding of the current state and potential evolution of graph prompts in AGI.

\subsection{Methodological Rigor}
Our survey methodology is grounded in a rigorous and systematic approach, ensuring the comprehensiveness and relevance of our analysis. We employed stringent criteria for literature selection, focusing on the scientific rigor, novelty, and impact of the included studies. This methodological rigor underpins the robustness and validity of our findings and conclusions.

% \section{Conclusion}
% This survey methodology section lays the groundwork for a detailed exploration of graph prompts within AGI. By outlining our research objectives, presenting a nuanced taxonomy, and providing a comprehensive literature overview, we set the stage for a profound understanding of the current landscape and future prospects in the field of graph prompts.